% Generated draft: requires mathematical and source review.
\documentclass[11pt]{article}
\usepackage[margin=25mm]{geometry}
\usepackage{fontspec}
\setmainfont{Latin Modern Roman}
\usepackage{amsmath,amssymb,mathtools,unicode-math}
\setmathfont{Latin Modern Math}
\usepackage{enumitem,needspace,xurl,hyperref}
\setlength{\parindent}{0pt}
\setlength{\parskip}{5pt}
\setlength{\emergencystretch}{3em}
\newcommand{\N}{\mathbb{N}}
\newcommand{\Z}{\mathbb{Z}}
\newcommand{\Q}{\mathbb{Q}}
\newcommand{\R}{\mathbb{R}}
\newcommand{\C}{\mathbb{C}}
\newcommand{\F}{\mathbb{F}}
\newcommand{\A}{\mathbb{A}}
\newcommand{\PP}{\mathbb{P}}
\newcommand{\E}{\mathbb{E}}
\newcommand{\eps}{\varepsilon}
\newcommand{\dd}{\,\mathrm{d}}
\begin{document}
\section*{UnsolvedMath ID: 30000792}
\textbf{Source problem code:} catalogue entry
\subsection{Definitions and mathematical statement}
Consider the source's space of finite-horizon dispersing billiard configurations in a flat $d$-torus. In a standard interpretation, particles move freely between smooth strictly convex scatterers and reflect specularly at collisions; finite horizon gives a uniform bound on free-flight length. The number, curvature and separation assumptions on scatterers and the precise configuration space are not supplied.

The question is whether configurations satisfying the report's finite-complexity condition form a generic subset in the $C^3$ topology. The source also gives the following as an equivalent formulation: a typical configuration should have no phase point whose trajectory undergoes more than $2d-2$ tangential collisions.

At a collision, tangency means that the incoming velocity is tangent to the scatterer boundary, equivalently its scalar product with the normal is zero. If $T(q,x)$ counts the intended tangential collision events for initial phase point $x$ in configuration $q$, the second formulation asks for generic $q$ satisfying
\[
 \forall x\text{ in the intended phase domain}:\qquad T(q,x)\le2d-2.
\]
The quantifier is every phase point in each generic configuration, not merely almost every phase point with respect to billiard measure. The time range, treatment of singular trajectories, continuation at a grazing collision, and counting convention must be recovered from the report.

The term generic might mean residual (containing a countable intersection of open dense sets), open dense, or another stated convention; the source does not specify which. The $C^3$ topology also requires a fixed parametrization or equivalence of obstacle boundaries. Neither a probability law on configurations nor a stronger open-dense assertion is inserted. The detailed finite-complexity condition is absent, so its asserted equivalence with the global tangency bound remains a supplied claim requiring review.

The dated assessment says the conjecture is proved for $d=2$ in the report but the multidimensional argument remains incomplete. This is retained as source context, without independently establishing the dimension-two result or narrowing the general question to a specific higher dimension.
\subsection{Short English statement}
The question asks whether generic finite-horizon dispersing billiards satisfy finite complexity, expressed by a uniform $2d-2$ bound on tangencies along every trajectory. The configuration space, meaning of generic, singular-trajectory rules, and exact finite-complexity equivalence are missing.
\subsection{Sources}
\begin{itemize}
\item Local UnsolvedMath registry, permanent ID 30000792, code catalogue entry, ``Generic Finite Complexity in Multidimensional Dispersing Billiards.'' Proposed year 2007; proposer unspecified.
\item \emph{Dynamische Systeme} (2008), Oberwolfach Reports, pp.\ 1913--1956, DOI 10.4171/owr/2007/33; \url{https://doi.org/10.4171/owr/2007/33}. The displayed report year 2008 differs from the DOI and proposed year 2007. This is the sole supplied literature reference.
\item Statement status corrected\_verified; the supplied note says official-report content was checked and wording/notation normalized without changing the source claim.
\item Literature triage checked 2026-08-21: partially solved. It reports a proof for $d=2$ in the cited report and explicitly incomplete higher-dimensional work. Exact hypotheses and current status require independent verification. Created 2026-08-25; updated 2026-09-16.
\item Registry curation by Ulam AI / UnsolvedMath Contributors, CC BY 4.0; the underlying report retains its supplied provenance. Unreviewed statement-only draft.
\end{itemize}
\end{document}
